\chapter{Introducción}

La transformación digital ha modificado sustancialmente la forma en que las organizaciones gestionan sus recursos tecnológicos. En las últimas décadas, la dependencia de sistemas informáticos ha crecido exponencialmente, convirtiendo a la tecnología en un pilar fundamental para la operación empresarial. Según Gartner, para 2025, el 85\% de las empresas tendría un principio cloud-first, priorizando el uso de tecnologías y servicios en la nube dentro de sus estrategias tecnológicas \parencite{Gartner2021}.

La computación en la nube surge como respuesta a la necesidad de flexibilidad y eficiencia en el acceso a recursos tecnológicos. Este paradigma permite provisionar infraestructura, plataformas y servicios bajo demanda, eliminando la necesidad de inversiones iniciales significativas en hardware físico. El National Institute of Standards and Technology (NIST) define cloud computing como ``un modelo para permitir acceso conveniente, bajo demanda y ubicuo a un pool compartido de recursos computacionales configurables que pueden ser rápidamente provisionados y liberados con mínimo esfuerzo de gestión o interacción con el proveedor de servicio'' \parencite{Mell2011}.

Las características fundamentales de la computación en la nube incluyen elasticidad, disponibilidad bajo demanda, pooling de recursos, medición de servicios y acceso amplio a través de la red. Estas propiedades permiten a las organizaciones escalar recursos vertical y horizontalmente según las necesidades operativas, optimizando tanto el rendimiento como los costos asociados. Sin embargo, esta flexibilidad introduce complejidades en la gestión que no estaban presentes en modelos tradicionales de infraestructura estática. Como señala Buyya et al., ``la naturaleza dinámica y elástica de las aplicaciones cloud y la provisionamiento de recursos en entornos heterogéneos y geográficamente distribuidos presenta desafíos significativos para el manejo eficiente de recursos'' \parencite{Buyya2013}.

La gestión de costos en entornos cloud representa uno de los principales desafíos operativos. El modelo de pago por uso, aunque económicamente ventajoso en teoría, puede resultar en gastos impredecibles sin mecanismos adecuados de monitoreo y control. Un estudio de Flexera revela que las organizaciones desperdician aproximadamente el 32\% de su gasto en cloud debido a recursos subutilizados y sobredimensionados \parencite{Flexera2023}. Esta situación evidencia la necesidad de herramientas que proporcionen visibilidad granular sobre el consumo de recursos y faciliten la identificación de ineficiencias.

La complejidad de las arquitecturas cloud modernas, caracterizadas por microservicios, contenedores y múltiples niveles de abstracción, dificulta la detección manual de anomalías en el consumo. Los responsables técnicos enfrentan el desafío de identificar instancias sobredimensionadas, recursos ociosos, configuraciones subóptimas y patrones de uso ineficientes. Según Amazon Web Services, ``la optimización continua de costos requiere un enfoque sistemático que combine monitoreo automatizado, análisis de datos y procesos de gobernanza'' \parencite{AWS2023}.

La automatización de infraestructura ha evolucionado como respuesta a estas complejidades. La práctica conocida como Infraestructura como Código (Infrastructure as Code - IaC) permite definir, desplegar y gestionar recursos tecnológicos mediante archivos de configuración versionables y reproducibles. Morris define IaC como ``la práctica de gestionar y provisionar recursos de infraestructura a través de archivos de configuración definidos por máquina, en lugar de herramientas interactivas o scripts manuales'' \parencite{Morris2016}. Terraform, desarrollado por HashiCorp, se ha consolidado como una de las herramientas más adoptadas para IaC, permitiendo gestionar infraestructura multi-cloud mediante un lenguaje declarativo.

El aprendizaje automático ofrece nuevas posibilidades para abordar la complejidad del análisis de infraestructura cloud. Las técnicas de machine learning, particularmente los algoritmos no supervisados, permiten identificar patrones anómalos en grandes volúmenes de datos operativos sin requerir etiquetado previo. El algoritmo Isolation Forest, propuesto por Liu et al. en 2008, detecta anomalías mediante el aislamiento de observaciones a través de particiones recursivas, resultando especialmente efectivo para detectar outliers en datasets multidimensionales \parencite{Liu2008}. Esta técnica resulta particularmente adecuada para identificar comportamientos anómalos en métricas de infraestructura, donde las ineficiencias se manifiestan como desviaciones respecto a patrones normales de operación.

En este contexto, el presente proyecto propone el diseño y desarrollo de una plataforma integrada para el monitoreo y optimización de recursos en entornos cloud. La solución combina tres pilares fundamentales: primero, la recolección sistemática de métricas operativas que proporcionen visibilidad sobre el estado y consumo de recursos; segundo, la aplicación del algoritmo Isolation Forest para detectar ineficiencias y comportamientos anómalos; y tercero, la generación de recomendaciones accionables orientadas a mejorar la utilización de recursos y optimizar costos. Adicionalmente, se incorpora una interfaz gráfica para visualización de datos y un módulo de despliegue automatizado basado en Terraform.

El desarrollo se enfoca en un prototipo funcional integrado con servicios cloud de Huawei, orientado a validar la viabilidad técnica de la propuesta en un entorno controlado. La selección de Huawei Cloud como plataforma obedece a que uno de los autores del presente trabajo se desempeña profesionalmente en dicha organización, lo cual proporciona acceso a créditos de prueba disponibles para la ejecución de las validaciones requeridas, así como conocimiento directo de los servicios y APIs de gestión de recursos del proveedor. El prototipo permitirá evaluar la efectividad de las técnicas propuestas y establecer las bases para desarrollos futuros que incorporen funcionalidades adicionales como optimización automática de recursos y predicción de costos.

La presente tesis se estructura en capítulos que abordan, en primer lugar, el marco conceptual necesario para comprender la problemática, incluyendo fundamentos de computación en la nube, gestión de costos y técnicas de aprendizaje automático. Posteriormente, se analizan las soluciones existentes en el mercado, identificando limitaciones y oportunidades de mejora. Los capítulos siguientes presentan el diseño arquitectónico de la plataforma, los detalles de implementación, las pruebas realizadas en entornos controlados y la evaluación de resultados. Finalmente, se discuten las conclusiones derivadas del trabajo y se proponen líneas de investigación futura.

Esta investigación contribuye al campo de la gestión de infraestructura cloud al proponer una solución que integra monitoreo, análisis inteligente mediante machine learning y automatización en una única plataforma. El enfoque adoptado busca demostrar que la combinación de estas tecnologías puede facilitar significativamente la administración de recursos cloud, promoviendo un uso más eficiente y contribuyendo a la reducción de costos operativos en organizaciones que adoptan modelos de computación en la nube.

\section{Objetivos}

\subsection{Objetivo General}

El objetivo general de este trabajo es diseñar y desarrollar una plataforma de optimización de infraestructura cloud que permita detectar recursos ineficientes mediante el algoritmo Isolation Forest, generar recomendaciones de optimización orientadas a reducir costos y mejorar la utilización de recursos, y proponer cambios de infraestructura mediante código de Terraform.

\subsection{Objetivos Específicos}

\begin{itemize}
\item Definir la arquitectura de la plataforma de optimización contemplando la recolección de métricas de consumo y rendimiento de recursos cloud.
\item Implementar mecanismos de recolección de datos sobre infraestructura desplegada en Huawei Cloud que permitan obtener información de utilización de recursos.
\item Desarrollar un módulo de análisis basado en el algoritmo Isolation Forest que permita detectar anomalías en el consumo de recursos, identificando instancias subutilizadas, sobredimensionadas o con patrones de uso ineficientes.
\item Implementar un componente de generación de recomendaciones que, a partir de las anomalías detectadas, proponga acciones concretas de optimización (rightsizing, eliminación de recursos ociosos, ajuste de configuraciones).
\item Desarrollar un módulo de integración con Terraform que permita generar código de infraestructura como código para implementar las recomendaciones de optimización propuestas.
\item Implementar una interfaz gráfica que permita visualizar el estado de la infraestructura, las anomalías detectadas, las recomendaciones generadas y el código Terraform propuesto.
\item Validar el funcionamiento del prototipo mediante pruebas en un entorno controlado, evaluando la capacidad del sistema para detectar ineficiencias y generar recomendaciones efectivas de optimización.
\end{itemize}
